% Part: proof-theory
% Chapter: normalization
% Section: translations

\documentclass[../../../include/open-logic-section]{subfiles}

\begin{document}

\olfileid{pt}{nor}{tr}

\olsection{सामान्यरूप \Log{N2i}-\usetoken{P}{proof} आणि \Log{G2i} यांतील रूपांतरे}

\Log{N2i} मधील !!^{proof}s चे \Log{G2i} मधील !!{proof}s मध्ये
आणि उलट दिशेने रूपांतर करता येते. प्रथम,
कट-विरहित \Log{G2i} !!{proof}s चे रूपांतर केल्यावर
सामान्यरूप \Log{N2i}-!!{proof}s मिळतात.

हा निष्कर्ष मांडण्यासाठी~\Log{N2i} आणि~\Log{G2i}
यांतील क्रमवर्तींचे पूर्वपक्ष जोडणारी काही परिभाषा हवी.
खूण लावलेल्या !!{formula}s चा संच~$\Gamma'$ हा खूण नसलेल्या
!!{formula}s च्या बहुसंच~$\Gamma$ शी \emph{अनुरूप} असतो,
जर प्रत्येक !!{formula}~$!A$ साठी~$\Gamma'$ मधील
$x : !A$ या रूपातील !!{element}s ची संख्या ही~$\Gamma$
मधील~$!A$ च्या प्रतींच्या संख्येहून जास्त नसेल.

\begin{cor}
$\Log{G2i} \Proves \Gamma \Sequent \Delta$ असेल, तर
$\Gamma' \Sequent !C$ ची सामान्यरूप
\Log{N2i}-!!{proof}~$\delta$ असते; येथे खूण लावलेल्या
!!{formula}s चा संच~$\Gamma'$ हा बहुसंच~$\Gamma$ शी अनुरूप असतो.
$\Delta=\emptyset$ असल्यास~$!C\ident\lfalse$,
आणि अन्यथा~$\Delta=\{!C\}$ असते.
म्हणजे, प्रत्येक !!{formula}~$!A$ साठी~$\Gamma'$ मधील
$x : !A$ या रूपातील !!{element}s ची संख्या ही~$\Gamma$
मधील~$!A$ च्या प्रतींच्या संख्येहून (म्हणजे~$!A$ च्या
पुनरावृत्तींच्या संख्येहून) जास्त नसते.
\end{cor}

\begin{proof}
\olref[nat][tgi]{prop:G2i-to-N2i} च्या सिद्धतेतील
रचनेकडे पाहिल्यास दिसते: !!{introduction} नियम
!!{proof}s च्या शेवटी, तर !!{elimination} नियम
फक्त !!{proof}s च्या सुरुवातीला जोडतो.
म्हणून !!a{elimination} नियमाच्या वर !!{introduction}
नियम कधी येत नाही.
\end{proof}

मात्र~\olref[nat][tgi]{prop:G2i-to-N2i} मधील
\Cut{} नियमाचे रूपांतर~\Log{N2i}-!!{proof}~$\delta$ मध्ये
कट निर्माण करू शकते. एखादी~\Log{G2i} !!{proof}
शेवटी~\Cut{} ने संपत असेल, आणि तिच्या तत्काळ उप-!!{proof}s
$\pi_1$ व~$\pi_2$ अनुक्रमे~$\Gamma_1 \Sequent !A$
आणि~$!A, \Gamma_2 \Sequent \Delta_2$ यांच्या असतील,
तर~$\pi_1$ चे रूपांतर~$\delta_1$ हे~$\pi_2$ च्या
रूपांतर~$\delta_2$ मधील $x:!A \Sequent !A$
स्वयंसिद्धकांवर कलम-संयोजित करतो. जर~$\delta_1$
मुख्य !!{formula}~$!A$ असलेल्या !!a{introduction}
नियमाने संपत असेल आणि~$\delta_2$ मधील
$x:!A \Sequent !A$ स्वयंसिद्धक हे !!a{elimination}
नियमाचे प्रमुख आधारविधान असेल, तर परिणामी~$\delta$
मध्ये कट तयार होतो.

\Log{N2i} मधील !!{proof}s चे~\Log{G2i} मधील
!!{proof}s मध्येही रूपांतर करता येते.
\olref[nat][tni]{prop:N2-to-G2} च्या सिद्धतेत
\Elim{\land}, \Elim{\lif}, \Elim{\lnot},
\Elim{\lforall} यांचे रूपांतर परिणामी~\Log{G2i}
!!{proof} मध्ये~\Cut{} अनुमाने निर्माण करते.
पण सुरुवातीची~\Log{N2i}-!!{proof} सामान्यरूप असेल,
तर असे कट टाळता येतात.

~\Log{N2i}-!!{proof}~$\delta$ मधील
क्रमवर्तींच्या आस्थितींच्या $S_1$, \dots,~$S_n$
या यादीला \emph{शाखा} म्हणू, जर~$S_1$ स्वयंसिद्धक
असेल आणि~$S_i$ जेव्हा एखाद्या अनुमानाचे प्रमुख
आधारविधान असेल तेव्हा~$S_{i+1}$ हा त्याचा निष्कर्ष
असेल. दुसऱ्या शब्दांत, शाखा स्वयंसिद्धकापासून
~$\delta$ मध्ये खाली जात क्रमवर्तींचा मागोवा घेते;
पण !!{elimination} नियमांच्या गौण आधारविधानांवर थांबते.
जिचा~$S_n$ अंतिम क्रमवर्ती असतो ती~$\delta$ ची
\emph{मुख्य शाखा}. प्रत्येक !!{proof}~$\delta$ ला
किमान एक मुख्य शाखा असते, कारण प्रत्येक नियमाला
किमान एक प्रमुख आधारविधान असते (फक्त~\Intro{\land}
ला दोन असतात).

\begin{prop}\ollabel{prop:normal-N2i-to-G2i}
~$\delta$ ही~$\Gamma \Sequent !A$ ची सामान्यरूप
$\Log{N2c}$ किंवा~$\Log{N2i}$-!!{proof} असेल,
तर~$\Gamma' \Sequent \Delta$ ची कट-विरहित
$\Log{G2c}$-!!{proof} (अनुक्रमे~$\Log{G2i}$-!!{proof})~$\pi$
असते. येथे~$\Gamma$ मधील खुणा काढून मिळणारा
!!{formula}s चा बहुसंच~$\Gamma'$ आहे;
$!A$ हा~$\lfalse$ नसेल तर~$\Delta = \{!A\}$,
आणि असेल तर~$\Delta = \emptyset$.
\end{prop}

\begin{proof}
यावेळी~$\Gamma \Sequent !A$ च्या !!{proof}~$\delta$
मधील अनुमानांच्या संख्येवर विगमन करतो.
~$\delta$ मध्ये अनुमान नसेल, तर ती
$x:!A \Sequent !A$ हा स्वयंसिद्धक आहे;
तेव्हा~$\pi$ हा अनुरूप~$!A \Sequent !A$
स्वयंसिद्धक आहे, किंवा~$!A \ident \lfalse$
असेल तर~$\lfalse \Sequent \ $ स्वयंसिद्धक आहे.

~$\delta$ चा शेवट अनुमाननियमात होत असेल,
तर पुढील प्रसंग वेगळे पाहू:
\begin{enumerate}
  \item $R$ हा !!a{introduction} नियम किंवा~\FalseInt
  असेल, तर रूपांतर~\olref[nat][tni]{prop:N2-to-G2}
  च्या सिद्धतेप्रमाणेच होते आणि~\Cut लागत नाही.

  \item ~$\delta$ चा शेवटचा नियम~$R$ हा !!a{elimination} नियम
  असेल, तर पुढील बाबी लक्षात घ्या:
  \begin{enumerate}
    \item ~$\delta$ ची मुख्य शाखा कोणत्याही
    !!{introduction} नियमातून जात नाही; अन्यथा त्या
    शाखेवर !!a{introduction} नियम किंवा~\FalseInt{}
    यानंतर !!a{elimination} नियम येईल आणि~$\delta$
    सामान्यरूप राहणार नाही.
    \item ~$\delta$ च्या मुख्य शाखांवर
    !!{introduction} नियम नसल्याने एकच मुख्य शाखा
    असते. (फक्त~\Intro{\land} ला दोन प्रमुख
    आधारविधाने आहेत; म्हणून अनेक मुख्य शाखा असतील,
    तर त्या~\Intro{\land} च्या आधारविधानांतून जाव्या लागतील.)
    \item मुख्य शाखेत~\Elim{\lor} किंवा~\Elim{\lexists}
    चे प्रमुख आधारविधान असेल, तर असा एकच नियम
    असतो आणि तो शेवटचा नियम~$R$ असतो. (नसल्यास
    शाखेत~\Elim{\lor} किंवा~\Elim{\lexists}
    चा निष्कर्ष पुढील !!a{elimination} नियमाचे प्रमुख
    आधारविधान होईल; क्रमपरिवर्तन रूपांतरण शक्य असल्यामुळे
    सिद्धता सामान्यरूप राहणार नाही.)
  \item \Elim{\lor} आणि~\Elim{\lexists} यांखेरीज
  कोणताही !!{elimination} नियम गृहीतके निरस्त करत नाही,
  म्हणजे क्रमवर्तींचा डावा संदर्भ बदलत नाही.
  \Elim{\lor} आणि~\Elim{\lexists} हे दोन्ही नियम
  फक्त गौण आधारविधानाच्या वरची
  गृहीतके निरस्त करतात; म्हणजे फक्त गौण
  आधारविधानांचे डावे संदर्भ बदलतात.
  त्यामुळे~$\delta$ च्या मुख्य शाखेवरील
  क्रमवर्ती~$S_i$ चा संदर्भ~$\Gamma_1$ असेल,
  तर~$S_i$ हे प्रमुख आधारविधान असलेल्या अनुमानाच्या
  निष्कर्ष~$S_{i+1}$ चा संदर्भ~$\Gamma_1' \supseteq \Gamma_1$
  असतो.
  \item म्हणून मुख्य शाखेतील सर्वांत वरचा
  क्रमवर्ती~$S_1$ हा~$x: !C \Sequent !C$
  असेल, तर~$x:!C \in \Gamma$.
  \end{enumerate}
  आता~$!C$ च्या रूपानुसार प्रसंग वेगळे पाहू.
  मुख्य शाखेवरील प्रत्येक~$S_i$ हे !!a{elimination}
  नियमाचे प्रमुख आधारविधान असल्याने
  $x:!C \Sequent !C$ बद्दलही तेच खरे आहे.

  \item $!C \ident !D \land !E$ असेल, तर~$\delta$
  चे रूप असे आहे:
  \[
  \Axiom$x:!D \land !E \fCenter !D \land !E$
  \RightLabel{\Elim{\land}[1]}
  \UnaryInf$x:!D \land !E \fCenter !D$
  \Deduce$x:!D \land !E, \Gamma_1 \fCenter !A$
  \DisplayProof
  \]
  मुख्य शाखेत~$x:!D \land !E \Sequent !D$ आहे.
  त्या शाखेत~\Intro{\lif}, \Intro{\lnot} यांचे
  प्रमुख आधारविधान किंवा~\Elim{\lor}, \Elim{\lexists}
  यांचे गौण आधारविधान येत नाही; म्हणून
  $x:!D \land !E$ मधील संदर्भ-!!{formula}s
  मुख्य शाखेवर बदलत नाहीत. त्यामुळे अंतिम
  क्रमवर्तीच्या संदर्भात~$x:!D \land !E$ असलेच पाहिजे.
  शिवाय,~\Elim{\land} ने संपणाऱ्या उप-!!{proof}
  ऐवजी~$x:!D \Sequent !D$ ठेवून, आणि मुख्य
  शाखेवरील प्रत्येक क्रमवर्तीच्या संदर्भात
  $x:!D \land !E$ ऐवजी~$x:!D$ ठेवून,
  !!{proof}~$\delta$ चे !!{proof}~$\delta_1$
  मध्ये रूपांतर करता येते; म्हणजे:
  \[
  \bottomAlignProof
  \Axiom$x:!D \land !E \fCenter !D \land !E$
  \RightLabel{\Elim{\land}[1]}
  \UnaryInf$x:!D \land !E \fCenter !D$
  \Deduce$x:!D \land !E, \Gamma_1 \fCenter !A$
  \DisplayProof
  \quad\text{ यात }\quad
  \bottomAlignProof
  \Axiom$x:!D \fCenter !D$
  \Deduce$x:!D, \Gamma_1 \fCenter !A$
  \DisplayProof
  \]
  विगमन गृहीतक~$\delta_1$ ला लागू होते,
  कारण~$\delta$ मधील अनुमानांची संख्या ही
  $\delta_1$ मधील अनुमानांच्या संख्येपेक्षा~$1$
  ने अधिक आहे.

  विगमन गृहीतकानुसार,
  $!D, \Gamma_1' \Sequent !A$ ची
  \Log{G2i}-!!{proof}~$\pi_1$ असते.
  \LeftR{\land} लावून~!!{proof}~$\pi$ मिळते:
  \[
  \AxiomC{}
  \RightLabel{$\pi_1$}
  \Deduce$!D, \Gamma_1' \fCenter !A$
  \RightLabel{\LeftR{\land}}
  \UnaryInf$!D \land !E, \Gamma_1' \fCenter !A$
  \DisplayProof
  \]
  $!A \ident \lfalse$ असेल, तर उदाहरणार्थ
  एक~\RightR{\Weakening} जोडतो:
  \[
  \AxiomC{}
  \RightLabel{$\pi_1$}
  \Deduce$!D, \Gamma_1' \fCenter $
  \RightLabel{\LeftR{\land}}
  \UnaryInf$!D \land !E, \Gamma_1' \fCenter $
  \RightLabel{\RightR{\Weakening}}
  \UnaryInf$!D \land !E, \Gamma_1' \fCenter !A$
  \DisplayProof
  \]
  \item $!C \ident !D \lor !E$ असेल,
  तर ते~\Elim{\lor} चे प्रमुख आधारविधान असले
  पाहिजे आणि हे अनुमान मुख्य शाखेवरील शेवटचे
  अनुमान~$R$ असले पाहिजे. म्हणजे~$\delta$ चे
  रूप असे आहे:
  \[
  \Axiom$x:!D \lor !E \fCenter !D \lor !E$
  \AxiomC{}
  \RightLabel{$\delta_1$}
  \Deduce$y:!D, \Gamma_1 \fCenter!A$
  \AxiomC{}
  \RightLabel{$\delta_2$}
  \Deduce$z:!E, \Gamma_2 \fCenter !A$
  \RightLabel{\Elim{\lor}}
  \TrinaryInf$x:!D \lor !E, \Gamma_1, \Gamma_2 \fCenter !A$
  \DisplayProof
  \]
  विगमन गृहीतक~!!{proof}s $\delta_1$ आणि~$\delta_2$
  यांना लागू होते; त्यामुळे अनुक्रमे
  $!D, \Gamma_1' \Sequent !A$ आणि
  $!E, \Gamma_2' \Sequent !A$ यांच्या
  \Log{G2i}-!!{proof}s $\pi_1$ व~$\pi_2$ मिळतात.
  त्यांच्यावर~$\LeftR{\lor}$ लावून~$\pi$ मिळते:
  \[
  \AxiomC{}
  \RightLabel{$\pi_1$}
  \Deduce$!D, \Gamma_1' \fCenter !A$
  \AxiomC{}
  \RightLabel{$\pi_2$}
  \Deduce$!E, \Gamma_2' \fCenter !A$
  \RightLabel{\LeftR{\lor}}
  \BinaryInf$!D \lor !E, \Gamma_1', \Gamma_2' \fCenter !A$
  \DisplayProof
  \]
  $\Elim{\lor}$ ने~$y: !D$ किंवा~$z: !E$
  निरस्त केले नसेल (म्हणजे ही गृहीतके गौण
  आधारविधानांच्या संदर्भात नसतील), किंवा~$!A$
  हे~$\lfalse$ असेल, तर काही~\LeftR{\Weakening}
  आणि~\RightR{\Weakening} अनुमाने जोडावी लागतात.
  उदाहरणार्थ:
  \[
  \AxiomC{}
  \RightLabel{$\pi_1$}
  \Deduce$\Gamma_1' \fCenter $
  \RightLabel{\LeftR{\Weakening}}
  \UnaryInf$!D, \Gamma_1' \fCenter $
  \AxiomC{}
  \RightLabel{$\pi_2$}
  \Deduce$\Gamma_2' \fCenter $
  \RightLabel{\LeftR{\Weakening}}
  \UnaryInf$!E, \Gamma_2' \fCenter $
  \RightLabel{\LeftR{\lor}}
  \BinaryInf$!D \lor !E, \Gamma_1', \Gamma_2' \fCenter $
  \RightLabel{\RightR{\Weakening}}
  \UnaryInf$!D \lor !E, \Gamma_1', \Gamma_2' \fCenter !A$
  \DisplayProof
  \]

  \item $!C \ident !D \lif !E$ असेल,
  तर~$\delta$ चे रूप असे आहे:
  \[
  \Axiom$x:!D \lif !E \fCenter !D \lif !E$
  \AxiomC{}
  \RightLabel{$\delta_1$}
  \Deduce$\Gamma_1 \fCenter!D$
  \RightLabel{\Elim{\lif}}
  \BinaryInf$x:!D \lif !E, \Gamma_1 \fCenter !E$
  \RightLabel{$\delta_2$}
  \Deduce$x:!D \lif !E, \Gamma_1, \Gamma_2 \fCenter !A$
  \DisplayProof
  \]
  येथे~$x:!D \lif !E, \Gamma_1 \Sequent !E$
  समाविष्ट करणाऱ्या मुख्य शाखेत~\Intro{\lif},
  \Intro{\lnot} यांचे प्रमुख आधारविधान किंवा
  \Elim{\lor}, \Elim{\lexists} यांचे गौण
  आधारविधान नसते. त्यामुळे~$x:!D \lif !E, \Gamma_1$
  मधील संदर्भ-!!{formula}s मुख्य शाखेवर बदलत
  नाहीत. म्हणून अंतिम क्रमवर्तीच्या संदर्भात
  $x:!D \lif !E, \Gamma_1$ असले पाहिजे.
  शिवाय,~\Elim{\lif} ने संपणाऱ्या उप-!!{proof}
  ऐवजी~$x:!E \Sequent !E$ ठेवून आणि मुख्य
  शाखेवरील प्रत्येक क्रमवर्तीच्या संदर्भात
  $x:!D \lif !E, \Gamma_1$ ऐवजी~$x:!E$
  ठेवून,~!!{proof} च्या~$\delta_2$ या भागाचे
  योग्य !!{proof}~$\delta_2'$ मध्ये रूपांतर
  करता येते; म्हणजे:
  \[
  \bottomAlignProof
  \Axiom$x:!D \lif !E, \Gamma_1 \fCenter !E$
  \RightLabel{$\delta_2$}
  \Deduce$x:!D \lif !E, \Gamma_1, \Gamma_2 \fCenter !A$
  \DisplayProof
  \quad\text{ यात }\quad
  \bottomAlignProof
  \Axiom$x:!E \fCenter !E$
  \RightLabel{$\delta_2'$}
  \Deduce$x:!E, \Gamma_2 \fCenter !A$
  \DisplayProof
  \]
  विगमन गृहीतक~$\delta_1$ आणि~$\delta_2'$
  यांना लागू होते; कारण~$\delta$ मधील
  अनुमानांची संख्या ही~$\delta_1$ व~$\delta_2$
  मधील अनुमानांच्या संख्यांची बेरीज अधिक~$1$
  इतकी आहे. हा एक अधिकचा टप्पा मुख्य शाखेच्या सुरुवातीच्या
  \Elim{\lif} अनुमानासाठी आहे. (तेच अनुमान
  $\delta$ मधील शेवटचे अनुमान~$R$ असेल,
  तर~$\delta_2$ मध्ये~$0$ अनुमाने असतात.)

  विगमन गृहीतकानुसार,
  \Log{G2i}-!!{proof}s मिळतात:
  $\Gamma_1' \Sequent !D$ ची~$\pi_1$
  आणि~$!E, \Gamma_2' \Sequent !A$ ची~$\pi_2$.
  \LeftR{\lif} लावून~!!{proof}~$\pi$ मिळते:
  \[
  \AxiomC{}
  \RightLabel{$\pi_1$}
  \Deduce$\Gamma_1' \fCenter !D$
  \AxiomC{}
  \RightLabel{$\pi_2$}
  \Deduce$!E, \Gamma_2' \fCenter !A$
  \RightLabel{\LeftR{\lif}}
  \BinaryInf$!D \lif !E, \Gamma_1', \Gamma_2' \fCenter !A$
  \DisplayProof
  \]
  $!D$ आणि/किंवा~$!A \ident \lfalse$ असेल,
  तर एक किंवा दोन~\RightR{\Weakening}
  अनुमाने जोडतो; उदाहरणार्थ:
  \[
  \AxiomC{}
  \RightLabel{$\pi_1$}
  \Deduce$\Gamma_1' \fCenter $
  \RightLabel{\RightR{\Weakening}}
  \UnaryInf$\Gamma_1' \fCenter !D$
  \AxiomC{}
  \RightLabel{$\pi_2$}
  \Deduce$!E, \Gamma_2' \fCenter $
  \RightLabel{\RightR{\Weakening}}
  \UnaryInf$!E, \Gamma_2' \fCenter !A$
  \RightLabel{\LeftR{\lif}}
  \BinaryInf$!D \lif !E, \Gamma_1', \Gamma_2' \fCenter !A$
  \DisplayProof
  \]

\end{enumerate}
उरलेले प्रसंग याच प्रकारे हाताळता येतात;
ते सरावासाठी ठेवले आहेत.
\end{proof}

\begin{cor}
सामान्यरूप~\Log{N1i} किंवा~\Log{N2i}-!!{proof}
मध्ये येणारे प्रत्येक !!{formula} हे अंतिम
!!{formula} किंवा अंतिम क्रमवर्ती यांचे उप-!!{formula}
असते.
\end{cor}

\begin{proof}
  \olref{prop:normal-N2i-to-G2i} नुसार प्रत्येक
  सामान्यरूप~\Log{N2i}-!!{proof} चे~\Log{G2i}-!!{proof}
  मध्ये रूपांतर करता येते. सिद्धतेकडे पाहिल्यास
  \Log{N2i}-!!{proof} मध्ये येणारे प्रत्येक !!{formula}
  \Log{G2i}-!!{proof} मध्येही येते. \Log{G2i}
  मध्ये~\Cut{} नियम नसल्याने त्याला उप-!!{formula}
  गुणधर्म आहे.
\end{proof}

\end{document}
